% !TeX root = ../main.tex
\section{Discussion}

The controlled verification establishes an executable link among declared architecture, supported observations, and change-time decisions. D1 checks whether graph and rule contracts agree on isolated patches. D2 reveals method-name lookalikes and unsupported Redis operations that ordinary topology patches may miss. P3 checks whether the incremental gate preserves this implementation's full-scan result. These checks answer different engineering questions; their shared implementation and development data prevent treating them as independent confirmations of general accuracy.

The relevant denominators are five changed nodes, 12 changed edges, seven violation rule sets, nine call-site cases and two anchor cases, and 20 positive plus 20 hard-negative detector signals. Repeated baseline structure inflates full-graph totals without increasing the diversity of changed relationships. D2 was deliberately designed from v0.1 weaknesses and informed v0.2. Its value is targeted diagnosis and regression protection. It is not a competitor baseline. A stronger comparison needs independent labels and a reviewed common observation unit; deterministic replay can reproduce the same wrong inference.

The relation-level gate permits gradual adoption: pre-existing keys remain visible without blocking unrelated changes. That choice also means an additional violating occurrence on the same edge may pass. The policy is therefore appropriate only when teams intend to gate new relationship-level findings, not every new violating call. An unresolved endpoint can escape observation altogether. Table~\ref{tab:policy-boundaries} makes these limits explicit so PASS cannot be mistaken for a complete architectural safety guarantee.

Separating violation from evolution avoids automatically forbidding every new infrastructure component. It does not automate architecture authority. The measured command neither authenticates a model-change approval nor records a durable rejection. A human rejection leaves unchanged inputs at REVIEW; repository controls must prevent the merge while code or contract changes are resolved. Later governance modules and protected branches cannot retrospectively establish the effectiveness of the measured workflow.

The observed 57/189 parsed-file ratio describes input scope, not CPU savings. The 53.2\% cold-to-warm nearest-rank p50 reduction combines parsing, Git startup, cache I/O, comparison, and reporting. Sequential observations on one small system do not isolate those costs or establish scalability.

External comparison remains future work. The next study must freeze a multi-repository holdout independent of detector tuning, blinded labels, tool versions and a non-empty capability-matched comparator subset. Report per-repository outcomes, unsupported cases, failures and uncertainty appropriate to sampling. A later practitioner pilot should measure false-block burden, evidence comprehension and review effort. Additional regression runs or citations cannot substitute for either study.
